\begin{textsample}{POS Dim 2 – persona\_grok – Score 25.00 – t195\_grok.txt}  \label{ex:f2_pos_040}
In the Philippines, \textbf{climate} change is upending the traditional \textbf{weather} \textbf{knowledge} that farmers and fishers have relied on for \textbf{generations}.
This \textbf{disruption} is putting food \textbf{security} and overall well-being at \textbf{risk}, as \textbf{communities} \textbf{face} more frequent typhoons, \textbf{floods}, and \textbf{droughts}.
These \textbf{challenges} are made worse by widespread \textbf{poverty} and inadequate governance, leaving \textbf{people} even more vulnerable.
The concept of"\textbf{resilience}"often gets thrown around in discussions about adapting to these \textbf{impacts}, but it places an unfair burden on affected \textbf{communities}.
It lets fossil \textbf{fuel} companies and \textbf{governments} off the hook, shifting responsibility away from those who have contributed most to the \textbf{crisis}.
Back in 2015, a petition was filed against major carbon-emitting companies, which paved the way for a groundbreaking \textbf{human} \textbf{rights} investigation in 2019.
This effort underscores the \textbf{push} for \textbf{accountability} in the \textbf{face} of \textbf{climate} harms.
The latest IPCC report echoes these concerns, warning of intensifying \textbf{climate} \textbf{risks} and increased mortality rates in vulnerable regions.
It emphasizes the urgent need for protecting \textbf{ecosystems}, implementing adaptation strategies, addressing \textbf{loss} and damage, and pursuing \textbf{climate} \textbf{justice} to keep global warming within 1.5°C.

% matched lemmas: accountability, challenge, climate, community, crisis, disruption, drought, ecosystem, face, flood, fuel, generation, government, human, impact, justice, knowledge, loss, people, poverty, push, resilience, right, risk, security, weather
\end{textsample}
